% concatenated from subgroup-tpp-ratio-of-nilpotent-groups-of-class-2.tex
%% filename: amsart-template.tex
%% version: 1.1
%% date: 2014/07/24
%%
%% American Mathematical Society
%% Technical Support
%% Publications Technical Group
%% 201 Charles Street
%% Providence, RI 02904
%% USA
%% tel: (401) 455-4080
%%      (800) 321-4267 (USA and Canada only)
%% fax: (401) 331-3842
%% email: tech-support@ams.org
%% 
%% Copyright 2008-2010, 2014 American Mathematical Society.
%% 
%% This work may be distributed and/or modified under the
%% conditions of the LaTeX Project Public License, either version 1.3c
%% of this license or (at your option) any later version.
%% The latest version of this license is in
%%   http://www.latex-project.org/lppl.txt
%% and version 1.3c or later is part of all distributions of LaTeX
%% version 2005/12/01 or later.
%% 
%% This work has the LPPL maintenance status `maintained'.
%% 
%% The Current Maintainer of this work is the American Mathematical
%% Society.
%%
%% ====================================================================

%     AMS-LaTeX v.2 template for use with amsart
%
%     Remove any commented or uncommented macros you do not use.

\documentclass{amsart}
\usepackage{amsmath, amssymb}
\usepackage[numbers]{natbib}
\usepackage{hyperref}
\usepackage{enumitem}
\usepackage{multirow}
\usepackage{float}
\usepackage{booktabs}
\usepackage{bm}

\setlist{nosep}
\bibliographystyle{authordate1}
\newtheorem{theorem}{Theorem}[section]
\theoremstyle{observation}
\newtheorem{observation}[theorem]{Observation}
\newtheorem{proposition}[theorem]{Proposition}
\newtheorem{lemma}[theorem]{Lemma}
\newtheorem{corollary}[theorem]{Corollary}
\theoremstyle{definition}
\newtheorem{definition}[theorem]{Definition}
\newtheorem{example}[theorem]{Example}
\newtheorem{conjecture}[theorem]{Conjecture}
\newtheorem{xca}[theorem]{Exercise}
\theoremstyle{remark}
\newtheorem{remark}[theorem]{Remark}
\numberwithin{equation}{section}

\begin{document}
\title[TPP Triples of Subgroups and Finite Nilpotent Groups of Class $2$]{On the triple product property for subgroups of finite nilpotent groups of class $2$}
\author{Sandeep R. Murthy}
\address{}
\curraddr{}
\email{srm@tuta.com}
\subjclass[2020]{Primary 20D60, Secondary 68R05}
\keywords{triple product property, finite groups, TPP triple, subgroup TPP triple, TPP ratio, subgroup TPP ratio, nilpotent groups, nilpotency class $2$, $p$-groups of nilpotency class $2$}
\date{16 January 2026}
\dedicatory{}

\begin{abstract}
{A number of upper bounds are proved relating to the triple product property (TPP) for subgroups of finite nilpotent groups of class $2$. The TPP is the property defined for three non-empty subsets $S, T, U$ of a group $G$ that the group equation $s's^{-1}t't^{-1}u'u^{-1} = 1$, over pairs of elements $s', s \in S$, $t', t \in T$, $u', u \in U$, is satisfied if and only if $s' = s$, $t' = t$, $u' = u$. When $G$ is finite, and the parameter $\rho_0(G)$, called \emph{subgroup TPP ratio}, is defined as $\rho_0(G) := \max\frac{|S||T||U|}{|G|}$, where the maximum is over the collection of all triples of subgroups $S, T, U$ of $G$ satisfying the TPP, this paper proves that \textup{(1)} $\rho_0(G) < \sqrt{|G:Z(G)}$} for (all) groups of nilpotency class $2$, \textup{(2)} $\rho_0(G) \leq p$ for $p$-groups with a cyclic commutator subgroup of order $p$, \textup{(3)} $\rho_0(G) = 1$ for $p$-groups of nilpotency class $2$ with a "large" centre, loosely defined as those satisfying $p^2 \leq |G:Z(G)| \leq p^3$, \textup{(4)} and $\rho_0(G) = 1$ for $p$-groups of nilpotency class $2$ with "small" (irreducible, complex) character degrees of $1$ or $p$.
\end{abstract}

\maketitle

\tableofcontents

\section{Introduction}

This paper addresses the following question(s):

\begin{quote}\label{Q1} If $G$ is \textup{(I)} a (finite) group of nilpotency class $2$, including, in particular, \textup{(II)} a $p$-group with a cyclic commutator subgroup of order $p$, or \textup{(III)} a $p$-group of nilpotency class $2$ with a "large" centre ($p^2 \leq |G:Z(G)| \leq p^3$), or \textup{(IV)} a $p$-group of nilpotency class $2$ with "small" character degrees ($\text{cd}(G) = \{1, p\}$), how large can the ratio $\rho_0(G) := \max \frac{|S||T||U|}{|G|}$ be, where the maximum is over the collection of all triples of subgroups $S, T, U$ of $G$ satisfying the so-called triple product property (TPP)? 
\end{quote}

These were motivated by patterns in computationally-derived TPP data for a number of "small" groups obtained originally by I. Hedtke and this author in 2012 \citep[Tables 1-4]{HM}, which included nilpotent groups of various nilpotency classes, as well as non-nilpotent groups. The data suggested the following propositions.

\textup{(1)} Groups with abelian normal subgroups of prime index $p$ satisfy $\rho_0 \leq \frac{p^2}{2p - 1}$, and, in particular, $p$-groups with abelian (maximal, normal) subgroups of index $p$ satisfy $\rho_0 = 1$.

\vspace{1mm}

\textup{(2)} Groups with cyclic normal subgroups of prime index $p$ satisfy $\rho \leq \frac{p^2}{2p - 1}$, where $\rho$ is a generalisation of $\rho_0$ to arbitrary subsets satisfying the TPP.

\vspace{1mm}

\textup{(3)} Groups of nilpotency class $2$ satisfy $\rho_0 < \sqrt{|G:Z(G)|}$ where $Z(G)$ is the group centre.

\vspace{1mm}

\textup{(4)} $p$-Groups with a cyclic commutator subgroup $G' = [G, G]$ of order $p$, which are necessarily of nilpotency class $2$, satisfy $\rho_0 \leq p$. This includes the extraspecial groups, which thus also satisfy $\rho_0 \leq p$.

\vspace{1mm}

\textup{(5)} $p$-Groups of nilpotency class $2$ with a large centre, say, such that $p^2 \leq |G:Z(G)| \leq p^3$, satisfy $\rho_0 = 1$. 

\vspace{1mm}

\textup{(6)} $p$-Groups of nilpotency class $2$ with small character degrees of $1$ or $p$ satisfy $\rho_0 = 1$.

\vspace{1mm}

Proposition \textup{(1)} was proved in \citep[Theorem 4.1, Corollary 4.3]{Mur}, while \textup{(2)} was conjectured for $p = 2$ in \citep[Conjecture 7.6]{HM} and for arbitrary $p$ in \citep[Conjecture 5.1]{Mur}. This paper proves propositions \textup{(3)-(6)} in sections \ref{SecGroupsOfTypeI}, \ref{SecGroupsOfTypeII}, \ref{SecGroupsOfTypeIII}, \ref{SecGroupsOfTypeIV}, with sections containing a table summarising the relevant computational data, where appropriate. There is a preliminary section containing definitions, basic facts and technical results needed for the main results.

\vspace{1mm}

The reader can refer to the sources for definitions and characterisations of nilpotency and nilpotency classes of groups, for example, the most standard one in terms of central series, and these are understood and taken as given. The focus of this paper is groups of nilpotency class $2$, which can be characterised as nonabelian groups satisfying the condition

\begin{equation}\label{NilpClassTwo}
\{1\} < G' \leq Z(G) < G.
\end{equation}

Particularly relevant for the approach used in this paper is that in such groups the centre $Z(G)$ is non-trivial, commutators are central, the central quotient $G/Z(G)$ is abelian, and there is an abelian quotient by any (normal) subgroup containing the centre. Bounds for the subgroup TPP ratio can then be derived by studying how the members of large subgroup TPP triples intersect with, or do not intersect with, the centre, and other key normal subgroups, such as the commutator subgroup $G'$.

\vspace{1mm}

Applications of the TPP for the complexity of matrix multiplication, for which it was originally introduced \citep{CU}, are not directly considered in the paper, but the upper bounds derived for the subgroup TPP ratio for the various classes of groups that are presented here are very relevant for such applications, for example, by leading to lower bounds for the parameter introduced in \cite[Definition 3.1, Lemma 3.1]{CU} for groups called the \emph{pseudoexponent} $\alpha$ that measures how effective a given group "realises" matrix multiplication. This paper continues the combinatorial and group-theoretic approach described in \citep{Neu2011}, and pursued in \citep{Mur}.

\section{Preliminaries}\label{Preliminaries}

Standard notation is used. Groups are generally finite, unless stated otherwise, and denoted by $G$, and subgroups by $H, K, N, P$ etc. or, in the case of subsets, including subgroups, satisfying the TPP, by $S, T, U$. The ring of integers modulo $n$ is denoted by $\mathbb{Z}_n$, and the (multiplicative) cyclic group of order $n$ by $C_n$. The term "$p$-group" denotes a finite group of order $p^n$ for some prime $p$ and positive integer $n$. All references to characters are to complex irreducible characters of finite groups. The set of unique character degrees of a group $G$ is denoted by $\text{cd}(G) = \{\chi(1) \mathrel | \chi \in \text{Irr}(G)\}$.

\vspace{1mm}

A number of basic definitions, notions and technical results used in the sections of main results will be recalled or proved, starting with the TPP \citep[Definition 2.1]{CU}.

\begin{definition}\label{TPP} Let $G$ be a group, finite or infinite, and $S, T, U$ non-empty subsets of $G$ with cardinalities $|S|, |T|, |U|$ respectively.  The triple $(S,T,U)$ satisfies the \emph{triple product property} (TPP) if there is an equivalence

\begin{equation}\label{TPPEqn} s's^{-1}t't^{-1}u'u^{-1} = 1 \iff s = s', t = t', u = u' \end{equation}

for all $s, s' \in S$, $t,t' \in T$, $u,u' \in U$. Then $G$ is said to \emph{realise} a TPP triple of \emph{parameter type}, or simply \emph{type}, $(|S|, |T|, |U|)$, and $|S|, |T|, |U|$ are called the \emph{parameters} of the triple and the product $|S||T||U|$ the \emph{size} of the triple.  If $S, T, U$ are subgroups of $G$ then $(S,T,U)$ is called a \emph{subgroup TPP triple} of $G$, for which the TPP condition simplifies to

\begin{equation} stu = 1 \iff s = t = u = 1, \end{equation}

for all $s \in S$, $t \in T$, $u \in U$. When $G$ is finite the TPP triple $(S, T, U)$ is be said to be of \emph{trivial} size or \emph{non-trivial} size depending on whether $|S||T||U| \leq |G|$ or $|S||T||U| > |G|$.
\end{definition}

Note that the term "subgroup TPP triple" refers only to TPP triples of subgroups, while "TPP triple" refers to triples composed of arbitrary (non-empty) subsets satisfying the TPP. Also note the permutation and translation invariance properties of TPP triples (see \citep[Lemma 2.1]{CU} and \citep[Observation 2.1]{Neu2011}) which are key in many results. As is a frequently used set-theoretic characterisation of the TPP reproduced below.

\begin{theorem} \label{HM} \textup{\citep[Theorem 3.1]{HM}} In a group $G$ three non-empty subsets $S,T,U \subseteq G$ satisfy the TPP if and only if 
\begin{equation}
\label{BasicTPPC2}Q(S) \cap Q(T)Q(U)  = Q(T) \cap Q(U) = \{1\}.
\end{equation}
where $Q(X) := \{x'x^{-1} \mathrel | x, x' \in X\}$ is called the quotient set of a non-empty subset $X\subseteq G$.  If $S, T, U$ are subgroups then this condition reduces to
\begin{equation}
\label{BasicTPPSub}S \cap TU  = T \cap U = \{1\}.
\end{equation}
\end{theorem}

A key consequence of \citep[Observation 2.1]{Neu2011} is that any TPP triple $(S', T', U')$ of a group $G$ can be translated to a type-equivalent TPP triple $(S, T, U)$ satisfying the property that

\begin{align}1 \in S \cap T \cap U.
\end{align}

Such TPP triples are called \emph{basic}. Subgroup TPP triples are necessarily basic. All references to TPP triples shall be to basic TPP triples, unless otherwise stated.

\vspace{1mm}

For a group $G$ the \emph{subgroup TPP ratio} is defined as

\begin{equation}\label{TPPRho0}
\rho_0(G) := \frac{\beta_0(G)}{|G|}
\end{equation}

where, following \citep{Neu2011, Mur}, $\beta_0(G)$ is the \emph{subgroup TPP capacity of $G$} defined as

\begin{equation}\label{TPPBeta0}
\beta_0(G) := \max \;\{|S||T||U| \mid (S,T,U) \textrm{ is a subgroup TPP triple of }G\}.
\end{equation}.

The second, more general quantity, called \emph{TPP ratio}, is defined as

\begin{equation}\label{TPPRho}
\rho(G) := \frac{\beta(G)}{|G|}
\end{equation}

where $\beta(G)$ is the \emph{TPP capacity of $G$} defined as

\begin{equation}\label{TPPBeta}
\beta(G) := \max \;\{|S||T||U| \mid (S,T,U) \textrm{ is a TPP triple of }G\}. 
\end{equation}

As $G$ always realises the trivial (subgroup) TPP triple $(G,\{1\},\{1\})$ then $|G| \leq \beta_0(G) \leq \beta(G)$ and $1 \leq \rho_0(G) \leq \rho(G)$. For abelian groups there is a basic but key fact.

\begin{theorem}\textup{\citep[from the proof of Lemma 3.1]{CU}}\label{ThmRhoAbelianGroups} If $G$ is an abelian group then $\rho_0(G) = \rho(G) = 1$, or, equivalently, $\beta_0(G) = \beta(G) = |G|$.
\end{theorem}
\begin{proof} If $(S, T, U)$ is a TPP triple of an abelian group $G$ then the "triple product map" $S \times T \times U \longrightarrow G$ given by $(s, t, u) \longmapsto (s, t, u)$ is injective.
\end{proof}

A further definition is useful in differentiating TPP triples by combinations of parameter type and size.

\begin{definition}\label{ProperTPP} A TPP triple $(S, T, U)$ of a group $G$ is said to be \emph{proper} if $\min\{|S|, |T|, |U|\} > 1$, or \emph{improper} if $\min\{|S|, |T|, |U|\} = 1$, or \emph{non-trivial} if $|S||T||U| > |G|$, or \emph{trivial} if $|S||T||U| \leq |G|$, or \emph{maximal} if $\rho(G) = \frac{|S||T||U|}{|G|}$ or, in case $(S, T, U)$ is a maximal subgroup TPP triple, if $\rho_0(G) = \frac{|S||T||U|}{|G|}$.
\end{definition}

Obviously, the proper/improper and trivial/non-trivial distinctions between TPP triples are mutually exclusive, but there are also combinations, for example, a proper or improper trivial TPP triple, or a non-trivial or trivial maximal TPP triple. Note that a maximal TPP triple need not be unique or non-trivial, and a non-trivial TPP triple need not be unique or maximal. But an improper TPP triple is always trivial \citep[Proof of Lemma 3.1]{CU}.

\vspace{1mm}

The notion of permutability of subsets of groups is very relevant to TPP triples, and is formally defined below.

\begin{definition} Two non-empty subsets $X, Y$ of a group $G$ are said to \emph{permute} (in $G$), or, equivalently, be \emph{permutable} (in $G$), if $XY = YX$, where $XY = \{xy \mathrel | x \in X, y \in Y\}$ and $YX = \{yx \mathrel | x \in X, y \in Y\}$. Clearly, this notion is symmetric, as $X$ permuting with $Y$ implies that $Y$ permutes with $X$.
\end{definition}

One form of permutability is when a subset normalises another: that is, for non-empty subsets $X, Y \subseteq G$ if $Xy = yX$ for all $y \in Y$ then $Y \subseteq N_G(X)$, where $N_G(X)$ is the normaliser of $X$ in $G$, and $XY = \bigcup_{y \in Y}Xy = \bigcup_{y \in Y}yX = YX$. Of course, if $X$ is a normal subgroup of $G$ then $N_G(X) = G$. A stronger form of permutability in groups, which implies normalisation, is centralisation, where, using the sets $X, Y$ above, $xy = yx$ for all $x \in X$ and $y \in Y$, in which case $Y \subseteq C_G(X)$, where $C_G(X)$ is the centraliser of $X$ in $G$ and $C_G(X) \leq N_G(X)$.

\begin{proposition}\label{PropTPPPermSubsets}Let $(S, T, U)$ be a non-trivial TPP triple of a group $G$, that is, one such that $|S||T||U| > |G|$. Then

\vspace{1mm}

\textup{(1)} No member $S, T, U$ permutes with another member and its inverse, that is, for any distinct $X, Y \in \{S, T, U\}$ either $XY \neq YX$ or $XY^{-1} \neq Y^{-1}X$.

\vspace{1mm}

\textup{(2)} No member $S, T, U$ is normalised by (or, in particular, centralised by) by $G$ or by each other, that is,  $N_G(X) < G$ for any $X \in \{S, T, U\}$, and $Y, Z \subsetneq N_G(X)$ for any distinct $X, Y, Z \in \{S, T, U\}$. In particular, if $(S, T, U)$ is a subgroup TPP triple then $S, T, U$ are all non-normal subgroups.

\vspace{1mm}

\textup{(3)} If $H$ is the (proper) normaliser (in $G$) of any member of $\{S, T, U\}$ and has index $n = |G: H|$ then $|S|$, $|T|$, $|U|$ satisfy
\begin{equation}
\frac{|S||T||U|}{|G|} \leq n\rho(H)
\end{equation}
and, furthermore, if $H$ is abelian, then
\begin{equation}
\frac{|S||T||U|}{|G|} \leq n.
\end{equation}
\end{proposition}
\begin{proof}Let $G$ and the TPP triple $(S, T, U)$ be as given.

\textup{(1)} Let $S$ permute with $T$ and $T^{-1}$. Then $Q(ST) = ST(ST)^{-1} = STT^{-1}S^{-1} = TT^{-1}SS^{-1} = Q(T)Q(S)$, where $Q(T)Q(S) \cap Q(U) = \{1\}$ (by the TPP for the permuted triple $(T, S, U)$), and so by \citep[Lemma 3.1]{CU} $(ST, U, \{1\})$ is a TPP triple of $G$ such that $|S||T||U| \leq |G|$, a contradiction. Thus $S$ does not permute with $T$ and $T^{-1}$, and, analogously, $S$ does not permute with $U$ and $U^{-1}$, and $T$ does not permute with $U$ and $U^{-1}$.

\vspace{1mm}

\textup{(2)} Let $S$ be normalised by $G$, that is, $N_G(S) = G$. Then $S$ permutes with $T$ and $T^{-1}$, and by part \textup{(1)} $|S||T||U| \leq |G|$, a contradiction. So $N_G(S) < G$, that is, $S$ is normalised by a proper subgroup of $G$. Analogously, $T$ and $U$ are not normalised by $G$, and $N_G(T), N_G(U) < G$. This means that if $S, T, U$ are subgroups they are all non-normal in $G$. In particular, $S$, $T$, and $U$ are not centralised by $G$, or by each other.

\vspace{1mm}

\textup{(3)} Let $H = N_G(S)$ and $n = |G:H|$. By part \textup{(1)}, $N_G(S) < G$, and also $T, U \subsetneq H$, that is, $T$ and $U$ are not contained in $H$ (otherwise $S$ would permute with $T$ and $T^{-1}$, or with $U$ and $U^{-1}$, and by part \textup{(1)} $|S||T||U| \leq |G|$, a contradiction).  Let $G/H$ be the collection of (left) cosets of $H$ in $G$, and let $\overline{S}, \overline{T}, \overline{U} \subseteq G/H$ be the $H$-supports of $S, T, U$ respectively, defined as in \citep[Definition 3.1]{Mur}, with sizes $\sigma = |\overline{S}|$, $\tau = |\overline{T}|$, $\tau = |\overline{U}|$ respectively. Now, $\sigma = 1$ as $S \subseteq H$, while $1 < \tau, \upsilon \leq n$ as $T, U \subsetneq H$. Also, let $T_0 = T \cap H$ and $U_0 = U \cap H$. Using independent right-translation of $T$ and $U$ \citep[Observation 4.1]{Neu2011} it can be supposed that $|T_0| \geq |T_y = T \cap yH|$ and $|U_0| \geq |U_y = U \cap zH|$ for any $yH \in \overline{T}$ and $zH \in \overline{U}$. The sets $\{T_y\}_{yH \in \overline{T}}$ and $\{U_z\}_{zH \in \overline{U}}$ form partitions of $T$ and $U$ respectively, and there is a bound for $|S||T||U|$ given by
\begin{align*}|S||T||U| &= |S|\sum_{yH \in \overline{T}}|T_y|\sum_{zH \in \overline{U}}|U_z| \\
                        &\leq \tau\upsilon|S||T_0||U_0| \\
                        &\leq n^2\beta(H).
\end{align*}

Dividing both sides by $|G| = n|H|$ yields
\begin{align*}\frac{|S||T||U|}{|G|} &\leq \frac{n^2\beta(H)}{|G|} \\
                                    &= \frac{n^2\beta(H)}{n|H|} \\
                                    &= n\rho(H)
\end{align*}


where, additionally, if $H$ is abelian, then $\rho(H) = 1$ and the conclusion follows.
\end{proof}

When at least two members of a TPP triple of a group are contained in a single larger subgroup the TPP ratio of the group is bounded by the TPP ratio of the subgroup, as described in the following result.

\begin{proposition}\label{PropTPPSplit} Let $G$ be a group, $H$ a subgroup, and $(S, T, U)$ a TPP triple of $G$. If $S, T \subseteq H$ then $\frac{|S||T||U|}{|G|} \leq \rho(H)$. Additionally, if $(S, T, U)$ is maximal for $G$ then $\rho(G) \leq \rho(H)$, and, additionally, if $H$ is abelian, then $\rho(G) = 1$. There is a subgroup version of this where "TPP triple" is replaced by "subgroup TPP triple", $\rho(G)$ by $\rho_0(G)$, and $\rho(H)$ by $\rho_0(H)$.
\end{proposition}
\begin{proof}Let $G$, $H$ and the TPP triple $(S, T, U)$ be given as above. Let $G/H$ be the collection of $n = |G:H|$ left cosets of $H$ in $G$, and let $\overline{U} = \{gH \in G/H \mathrel | U \cap gH \neq \emptyset\}$ be the $H$-support of $U$. As $1 \in U \cap H$ then $U$ is at least partly contained in $H$ and $U_0 = U \cap H$ is non-empty. Then $(S, T, U_0)$ forms a TPP triple of $H$ such that $\frac{|S||T||U_0|}{|H|} \leq \rho(H)$, and, by translation invariance \citep[Observation 2.1]{Neu2011} it may be supposed that any support element $gH \in \overline{U}$ intersects with $U$ in at most $|U_0|$ elements, that is, $|U \cap gH| \leq |U_0|$ for all $gH \in \overline{U}$. The set $\overline{U}$ forms a partition of $U$ over the cosets of $H$ in $G$, and $|U| \leq n|U_0|$. Letting $U_1 = U \backslash U_0$ and $|U_1| = |U| - |U_0|$, so that $|U| = |U_0| + |U_1|$, where $|U_1| \leq (n - 1)|U_0|$, then

\begin{align*}\frac{|S||T||U|}{|G|} &= \frac{|S||T|(|U_0| + |U_1|)}{|G|} \\
                                    &= \frac{|S||T|(|U_0| + |U_1|)}{n|H|} \\
                                    &= \frac{1}{n}\frac{|S||T||U_0|}{|H|} + \frac{1}{n}\frac{|S||T||U_1|}{|H|} \\
                                    &\leq \frac{1}{n}\frac{|S||T||U_0|}{|H|} + \frac{n - 1}{n}\frac{|S||T||U_0|}{|H|} \\
                                    &= \frac{|S||T||U_0|}{|H|} \\
                                    &\leq \rho(H).
\end{align*}

The second part of the claim follows by taking $(S, T, U)$ to be maximal for $G$, that is, such that $\rho(G) = \frac{|S||T||U|}{|G|}$. From this the third part follows by taking $H$ to be abelian, in which case $\rho(H) = \rho(G) = 1$. The equivalent statement for subgroups follows by replacing "TPP triple" by "subgroup TPP triple", $\rho(G)$ by $\rho_0(G)$, and $\rho(H)$ by $\rho_0(H)$.
\end{proof}

There is a corollary - also independently implied by Proposition \ref{PropTPPPermSubsets} - that is not needed for the main results but is useful to note.

\begin{corollary} If $(S, T, U)$ is a non-trivial TPP triple of a nonabelian group $G$, that is, such that $|S||T||U| > |G|$, then no two members $X, Y \in \{S, T, U\}$ generate an abelian subgroup, that is, the subgroups $\langle S, T \rangle$, $\langle S, U \rangle$, $\langle T, U \rangle$ are all nonabelian.
\end{corollary}
\begin{proof} Use Proposition \ref{PropTPPSplit} and let $H = \langle X, Y \rangle$ for any two $X, Y \in \{S, T, U\}$.
\end{proof}

A basic fact about abelian normal subgroups in groups of nilpotency class $2$ realised as direct products that will be useful.

\begin{observation}\label{ObsNilpClassTwo} If $G$ is a group of nilpotency class $2$, and $H$ a subgroup of $G$ such that $H \cap Z(G) = \{1\}$, then $HZ(G)$ is an abelian normal subgroup of $G$ and a direct product.
\end{observation}
\begin{proof}Let $G$ and $H$ be given as above, and let $H \cap Z(G) = \{1\}$. As $G' \leq Z(G)$ by assumption, $H' = [H, H] \leq H \cap G' = \{1\}$ and $H$ is abelian. As $Z(G)$ is normal in $G$ the product $HZ(G)$ is a subgroup of $G$. The quotient $G/Z(G)$ is abelian, $HZ(G)/Z(G)$ is normal in $G/Z(G)$ and, therefore, $HZ(G)$ is normal in $G$ (by the Correspondence Theorem between the subgroups of $G$ containing $Z(G)$ and the subgroups of $G/Z(G)$). Finally, as $H \cap Z(G) = \{1\}$ then $HZ(G)$ is an abelian direct product of order $|HZ(G)| = |H||Z(G)|$.
\end{proof}

The following is a very useful result \citep{Neu2012} that is a kind of conditional converse to \citep[Lemma 2.2]{CU}.

\begin{lemma}\textup{\citep{Neu2012}}\label{LemNeuQuoSubTPP} Let $(S, T, U)$ be a subgroup TPP triple of $G$ and $N$ a normal subgroup of $G$ such that $N \subseteq S$. Then the quotient $G/N$ realises the subgroup TPP triple $(S/N,TN/N,UN/N)$ of type $\left(|S|/|N|, |T|, |U|\right)$, and $\rho_0(G/N) \geq \frac{|S||T||U|}{|G|}$. In particular, if $(S, T, U)$ is maximal for $G$ then $\rho_0(G/N) \geq \rho_0(G)$.
\end{lemma}
\begin{proof} Recall that $S, T, U$, as TPP subgroups, satisfy $S \mathrel \cap TU = T \mathrel \cap U = \{1\}$. As $N \subseteq S$ by assumption then $TU \cap N =\{1\}$. 

\vspace{1mm}

Let $\pi:G \rightarrow G/N$ be the canonical projection of $G$ onto $G/N$, that is, $\pi(g) = gN$ for each $g \in G$, and let $\overline{S}, \overline{T}, \overline{U} \subseteq G/N$ be the images of $S, T, U$ respectively, under $\pi$, in $G/N$, that is, $\overline{S} = \pi(S) = \{sN \mathrel | s \in S\}$, and $\overline{T} = \pi(T) = \{tN \mathrel | t \in T\}$, and $\overline{U} = \pi(U) = \{uN \mathrel | u \in U \}$. Then $\overline{S} = SN/N = S/N$, and $\overline{T} = TN/N \cong T/\{1\} \cong T$, and $\overline{U} = UN/N \cong U/\{1\} \cong U$, and these are all subgroups of $G/N$.

\vspace{1mm}

Letting $1_{G/N} = N$ be the identity of $G/N$ then $\overline{T} \cap \overline{U} \cong T \cap U = \{1\}$, which implies that $\overline{T} \cap \overline{U} = \{1_{G/N}\}$. Consider $\overline{S} \cap \overline{T}\cdot \overline{U} = S/N \cap (TN/N \cdot UN/N) = S/N \cap (TU)N/N$, where $(TU)N/N = \{tuN \mathrel | t \in T, u \in U\}$ is a subset of $G/N$ and not necessarily a subgroup (unless $G/N$ is abelian). Let $xN \in S/N \cap (TU)N/N$. Then $xN = sN = tuN$ for some $s \in S$ and $tu \in TU$ for some $t \in T$ and $u \in U$, which implies that $tu \in sN \cap TU \subseteq S \cap TU = \{1\}$, that is, $tu = 1$, which, by the TPP for $(S, T, U)$, implies that $s = t = u = 1$. This shows that $S/N \cap (TU)N/N = (S \cap T)N/N = N/N \cong \{1\}$ and $\overline{S} \cap \overline{T}\cdot \overline{U} = \{1_{G/N}\}$, where $\overline{T}\cdot\overline{U} \cong TU$ is a set bijection, and that $(\overline{S}, \overline{T}, \overline{U})$ is a subgroup TPP triple of $G/N$ of type $(|S|/|N|, |T|, |U|)$ as claimed.

\vspace{1mm}

The last part of the claim follows by taking $(S, T, U)$ to be a maximal subgroup TPP triple of $G$.
\end{proof}

This has two simple corollaries.

\begin{corollary}\label{Corr1LemNeuQuoSubTPP} Any group $G$ contains a normal subgroup $N$ such that $\rho_0(G/N) \geq \rho_0(G)$.
\end{corollary}
\begin{proof} Let $(S, T, U)$ be a subgroup TPP triple of $G$ such that $\rho_0(G) = \frac{|S||T||U|}{|G|}$ and $N = \text{Core}_G(X) = \bigcap_{g \in G}g^{-1}Xg$ the normal $G$-core of any member $X \in \{S, T, U\}$, which is the largest normal subgroup of $G$ contained in $X$, formed by the intersection of all the conjugates of $X$ by elements of $G$: any member, $X = S$, say, may be chosen, in which case $N = \text{Core}_G(S)$ and by Lemma \ref{LemNeuQuoSubTPP} $G/N$ realises the subgroup TPP triple $(S/N,TN/N, UN/N)$ of size $(|S||T||U|)/|N|$, and $\rho_0(G/N) \geq \frac{(|S||T||U|)/|N|}{|G|/|N|} = \rho_0(G)$. For the tightest bound $X \in \{S, T, U\}$ can be chosen to maximise the order of $N = \text{Core}_G(X)$.
\end{proof}

Note above that if $N = \text{Core}_G(S) = \{1\}$ then $S$ is a core-free subgroup and $G$ is isomorphic to a transitive subgroup of $\text{Sym}(n)$ where $n = |G:S|$.

\begin{corollary}\label{Corr2LemNeuQuoSubTPP} Let $G$ be a nonabelian group and $(S, T, U)$ a non-trivial subgroup TPP triple of $G$, that is, one such that $|S||T||U| > |G|$. Then 

\vspace{1mm}

\textup{(1)} No member $S, T, U$ contains the commutator subgroup $G'$.

\vspace{1mm}

\textup{(2)} No member $S, T, U$ contains a normal subgroup $N$ of $G$ such that the quotient $G/N$ is abelian.

\vspace{1mm}

\textup{(3)} If $G$ is of nilpotency class $2$ then no member $S, T, U$ contains the centre $Z(G)$.

\vspace{1mm}

\textup{(4)} If $G$ is a $p$-group then no member $S, T, U$ contains the Frattini subgroup $\Phi(G)$.
\end{corollary}
\begin{proof} Let $G$ and $(S, T, U)$ be as given.

\vspace{1mm}

\textup{(1)} By Theorem \ref{ThmRhoAbelianGroups}) $\rho_0(H) = \rho(H) = 1$ if $H$ is abelian. The commutator subgroup $G'$ is the minimal normal subgroup of $G$ such that the quotient $G/G'$ is abelian, which means that $\rho_0(G/G') = 1$. Choosing any member, $S$, say, such that $G' \subseteq S$, by Lemma \ref{LemNeuQuoSubTPP} $G/G'$ realises the subgroup TPP triple $(S/G', TG'/G', UG'/G')$ of type $(|S|/|G'|, |T|, |U|)$, and $\frac{|S||T||U|}{|G|} = \frac{(|S||T||U|)/|G'|}{|G|/|G'|} \leq \rho_0(G/G') = 1$, a contradiction. So $S$, and therefore also $T$ and $U$, cannot contain $G'$, from the assumption $|S||T||U| > |G|$.

\vspace{1mm}

\textup{(2)} Follows from \textup{(1)}.

\vspace{1mm}

\textup{(3)} If $G$ is of nilpotency class $2$ then the centre $Z(G)$ contains the commutator subgroup $G'$ and the conclusion follows from \textup{(1)}.

\vspace{1mm}

\textup{(4)} If $G$ is a (nonabelian) $p$-group then the Frattini subgroup $\Phi(G)$ contains a non-trivial commutator subgroup $G'$, and the conclusion follows from \textup{(1)}.
\end{proof}

The following lemma relates the subgroup TPP ratio of a quotient of a group, by some normal subgroup, to the index of its centre, and will be useful.

\begin{lemma}\label{LemQuoNilpClassTwo} If $G$ is a group with a normal subgroup $N$ such that $\rho_0(G/N) \leq \sqrt{|G/N|:|Z(G/N)|}$ then $\rho_0(G/N) \leq \sqrt{|G:Z(G)|}$.
\end{lemma}
\begin{proof} Let $G$, $N$ and $\rho_0(G/N)$ be given as above. Note that $Z(G)N/N \leq Z(G/N)$, where $Z(G)N/N \cong Z(G)/(Z(G) \cap N)$ and $|Z(G)N/N| = \frac{|Z(G)|}{|Z(G) \cap N|}$, and therefore $\frac{1}{|Z(G/N)|} \leq \frac{1}{|Z(G)N/N|} = \frac{|Z(G) \cap N|}{|Z(G)|}$. Now
\begin{align*}\rho_0(G/N) &\leq \sqrt{|G/N|:|Z(G/N)|} \\
                          &= \frac{\sqrt{|G|}}{\sqrt{|N|}}\cdot\frac{1}{\sqrt{|Z(G/N)|}} \\
                          &\leq \frac{\sqrt{|G|}}{\sqrt{|N|}}\cdot\frac{1}{\sqrt{|Z(G)N/N|}} \\
                          &= \frac{\sqrt{|G|}}{\sqrt{|N|}}\cdot\frac{\sqrt{|Z(G) \cap N|}}{\sqrt{|Z(G)|}} \\
                          &\leq \frac{\sqrt{|G|}}{\sqrt{|N|}}\cdot\frac{\sqrt{|N|}}{\sqrt{|Z(G)|}} \\
                          &= \sqrt{|G:Z(G)|}.
\end{align*}
\end{proof}

Two commutator identities for groups of nilpotency class $2$ will be used, one following from the other. The first is

\begin{equation}\label{EqCommId} [xy, z] = [x, z][y, z], \hspace{3em}x, y, z \in G \text{ (where $G$ is of nilpotency class $2$)}
\end{equation}

This uses the fact that, in this case, commutators are central and for elements $x, y, z \in G$ the identity is given by $[xy, z] = (xy)^{-1}z^{-1}xyz = y^{-1}x^{-1}z^{-1}xyz = y^{-1}x^{-1}z^{-1}xzz^{-1}yz = y^{-1}[x, z]z^{-1}yz = [x, z]y^{-1}z^{-1}yz = [x, z][y, z]$. From this follows the second

\begin{equation}\label{EqCommId2} [x^p, y] = [x, y]^p.
\end{equation}

For the latter note that by \ref{EqCommId} $[x^p, y] = [x\cdot x^{p - 1}, y] = [x, y][x^{p - 1}, y]$, and repeated application leads to $[x, y]^p$, which can be proved by induction.

\vspace{1mm}

The following proposition is useful to note for some of the later results for $p$-groups.

\begin{proposition}\label{PropRho0p4} If $G$ is a $p$-group of order $\leq p^4$ then $\rho_0(G) = 1$.
\end{proposition}
\begin{proof} All $p$-groups of order $\leq p^2$ are abelian and realise $\rho_0 = \rho = 1$. For nonabelian $p$-groups of order $p^3$ and $p^4$ the result follows from combinatorial constraints on parameters of TPP triples imposed by the inequalities in \citep[Observation 3.1]{Neu2011} by which the possible parameter types for subgroup TPP triples of groups of these orders, where all members are non-trivial proper subgroups, are $(p, p, p)$ and $(p^2, p, p)$ respectively, in both cases implying $\rho_0 = 1$. Some simple group-theoretic arguments also exist here. Nonabelian $p$-groups of order $p^3$, which are extraspecial, and those of order $p^4$, contain abelian (and normal) subgroups of index $p$ (see \citep[Lemma 3.4]{Cra} for $p$-groups of order $p^3$, and for $p$-groups of order $p^4$ see \citep[Lemma 2]{Wil} for $p = 2$, and \citep[Proposition 12]{AGW} for $p > 2$). So for these groups $\rho_0 = 1$ by \citep[Corollary 4.3]{Mur}.
\end{proof}

For $p$-groups with a cyclic commutator subgroup of order $p$, or $p$-groups of nilpotency class $2$ that satisfy a large centre assumption $|G:Z(G)| \leq p^3$, some basic results are required, all of which are known or can be deduced from the basic theory of (finite) $p$-groups.

\begin{lemma}\label{LemCycNormSubOrderp} If $N$ is a normal subgroup of a nonabelian $p$-group $G$ such that $N$ is cyclic of order $p$ then $N$ is central, that is, $N \trianglelefteq G$ and $N \cong \mathbb{Z}_p$ implies that $N \leq Z(G)$. In particular, if $G$ is a nonabelian $p$-group containing a cyclic commutator subgroup $G'$ of order $p$ then $G'$ is central and $G$ is of nilpotency class $2$.
\end{lemma}
\begin{proof} Let $G$ and $N$ be as given. If $G$ is abelian $Z(G) = G$ and there is nothing to prove. So suppose $G$ is nonabelian. Then $G$ acts on the elements of $N$ by conjugation, given by $n^g \longmapsto g^{-1}ng = n'$ for elements $n \in N$ and $g \in G$, where $n' \in N$ depends on $n$, and produces conjugacy classes $n^G = \{g^{-1}ng \mathrel | g \in G\}$ for $n \in N$, that partition $N$. The number $k$ of these conjugacy classes satisfies $1 \leq k \leq p = |N|$, as there is a trivial singleton conjugacy class $1^G = \{1\}$, and each conjugacy class size $|n^G|$ must divide $|G|$ and is therefore either $1$ or $p$. This implies that every non-identity element of $N$ also forms a singleton conjugacy class, that is, all elements of $N$ are self-conjugate: $g^{-1}ng = n$ for any $n \in N$ and $g \in G$. This is equivalent to $[n, g] = n^{-1}g^{-1}ng = 1$ for all $g \in G$ and $n \in G$, that is, all elements of $N$ commute with all elements of $G$, formally, $[N, G] = \{1\}$, which proves that $N \leq Z(G)$. The second part of the claim follows by substituting $G'$ for $N$, in which case there is a central series for $G$ given by $1 < G' \leq Z(G) < G$ of length $2$.
\end{proof}

\begin{lemma}\label{LemNonAbCentQuo} If $G$ is a nonabelian group then $G/Z(G)$ cannot be cyclic.
\end{lemma}
\begin{proof}
Suppose $G$ is nonabelian. If $G/Z(G)$ is cyclic, say, of order $n = |G:Z(G)|$, then $G/Z(G) = \langle g_0Z(G) \rangle = \{(g_0Z(G))^i = g_0^iZ(G) \mathrel | i \in \mathbb{Z}_n\}$ for some $g_0 \in G\backslash\{1\}$, and $G = \bigcup_{i \in \mathbb{Z}_n}g_0^iZ(G)$, which means that any element in $G$ can be written as $g_0^iz$ for some $i \in \mathbb{Z}_n$ and $z \in Z(G)$ depending on the element. If $x, y \in G$ are any two elements then $x = g_0^iz_x$ and $y = g_0^jz_y$ for some $i, j \in \mathbb{Z}_n$ and $z_x, z_y \in Z(G)$, and $xy = g_0^iz_x \cdot g_0^jz_y = g_0^ig_0^jz_xz_y = g_0^{i + j}z_xz_y = g_0^jg_0^iz_yz_x = g_0^jz_y \cdot g_0^iz_x = yx$, which implies that $G$ is abelian, a contradiction.
\end{proof}

\begin{proposition}\label{PropAbCentQuop2p3} If $G$ is a $p$-group such that $G/Z(G)$ is elementary abelian of order $p^2$ or abelian of order $p^3$ then $G$ contains abelian (maximal, normal) subgroups of index $p$.
\end{proposition}
\begin{proof} Let $G$ be as given, in which case $G$ is of nilpotency class $2$. Let $G/Z(G)$ be elementary abelian of order $p^2$. Then $G/Z(G)$ is isomorphic to $C_p \times C_p$, and contains an abelian maximal normal subgroup, $L$, say, of index $p$ that is cyclic of order $p$ and generated by some element $xZ(G)$ of order $p$. By the Correspondence Theorem the subgroup of $G$ corresponding to $L$ is an abelian maximal normal subgroup $K$ of index $p$ that is generated by $Z(G)$ and some non-central element $x \in G\backslash Z(G)$ of order $\geq p$, where $K$ is such that $K = C_G(x)$ where $C_G(x)$ is the centraliser of $x$ in $G$. \textup{(2)}. Let $G/Z(G)$ be abelian of order $p^3$. Then $G/Z(G)$ is isomorphic to either $C_{p^2} \times C_p$ or $C_p \times C_p \times C_p$. If $G/Z(G) \cong C_{p^2} \times C_p$ then $G$ contains an abelian maximal normal subgroup $K$ of index $p$ corresponding to an abelian maximal normal subgroup $L$ of $G/Z(G)$ of index $p$ where $L$ is isomorphic to $C_{p^2}$, in which case $K$ is generated by $Z(G)$ and some non-central element $x \in G \backslash Z(G)$ of order $\geq p^2$, where $K = C_G(x)$. If $G/Z(G) \cong C_p \times C_p \times C_p$ then $L$ is isomorphic to $C_p \times C_p$, and $K$ is generated by $Z(G)$ and two non-central commuting elements $x, y \in G \backslash Z(G)$ of order $\geq p$, where $K = C_G(\{x, y\})$.
\end{proof}

For the case of extraspecial groups, a little lemma on a non-generating property of two "large" abelian subgroups will be useful.

\begin{lemma}\label{LemExspSub} Let $G$ be an extraspecial group, which is necessarily of order $p^{1 + 2n}$ for some $n \geq 1$. If $S$ and $T$ are two subgroups of $G$ of order $p^n$ then $S$ and $T$ do not generate $G$, that is, $\langle S, T \rangle < G$.
\end{lemma}
\begin{proof} Let $G$, $S$ and $T$ be as given. Then $G$ is an (internal) central product $P_1 * \cdots * P_n$ of $n$ extraspecial factor subgroups $P_1,\ldots,P_n$ each of order $p^3$, which have these properties: each is normal in $G$ as each contains $G'$, each satisfies $P_i' = Z(P_i) = \Phi(P_i) \cong \mathbb{Z}_p$ so that the quotient $G/P_i$ is (elementary) abelian, and any (distinct) pair of them centralise each other, that is, $[P_i, P_j] = \{1\}$ for $1 \leq i, j \leq n$ such that $i \neq j$.\footnote{The main references for central products, $p$-groups and extraspecial groups are \citep[Chapter 3]{Cra} and \citep[Chapter 4]{Ber}} It is also the case that $G$ is of nilpotency class $2$ since $\{1\} < G' = P'_i = Z(P_i) = Z(G) < G$, where $P_i$ is any of the extraspecial factors.

\vspace{1mm}

Suppose $S$ and $T$ generate $G$. As $G' = Z(G)$ the quotient $G/Z(G)$ is (elementary) abelian, and the products $SZ(G)$ and $SZ(G)T = STZ(G)$ are normal subgroups of $G$, and $G = \langle S, T \rangle = SZ(G)T$. The order of $G$ is $|G| = |SZ(G)T| = \frac{|SZ(G)||T|}{|SZ(G) \cap T|} = \frac{|S||Z(G)||T|}{|S \cap Z(G)||SZ(G) \cap T|} = \frac{p^{2n + 1}}{|S \cap Z(G)||SZ(G) \cap T|}$, which implies that $S \cap Z(G) = SZ(G) \cap T = \{1\}$, and in particular $S \cap T = \{1\}$. Transposing $S$ and $T$ implies $T \cap Z(G) = TZ(G) \cap S = \{1\}$. Also, $S$ and $T$ are abelian since, to take $S$ as an example, $S' = [S, S] \leq S \cap G' = S \cap Z(G) = \{1\}$. As $G$ is nonabelian and $G' = Z(G)$ is cyclic of order $p$ then $[S, T] > \{1\}$ and $[S, T] = G' = Z(G)$.

\vspace{1mm}

For an extraspecial factor, $P_1$, say, consider the subgroup $SP_1T = G$. Then $|G| = |SP_1T| = \frac{|S||P_1||T|}{|S \cap P_1||SP_1 \cap T|} = \frac{p^{2n + 3}}{|S \cap P_1||SP_1 \cap T|}$, which implies $|S \cap P_1||SP_1 \cap T| = p^2$ and the subsets $S \cap P_1$ and $SP_1 \cap T$ are non-empty. Suppose $S \cap P_1$ is non-empty. Then a non-identity element $x_1 \in S \cap P_1$ can be chosen, where $x_1 \not\in Z(G)$, such that $\{1\} = [x_1, P_2] = \cdots = [x_1, P_n]$ and $x_1 \in Z(\langle x_1, P_2, \ldots, P_n \rangle) \leq Z(P_1* \cdots * P_n) = Z(G)$, a contradiction. Thus, $S \cap P_1 = \{1\}$ and in the same way it can be shown that $T \cap P_1 = \{1\}$. Then $|SP_1 \cap T| = p^2$ and a non-identity element $y_1 \in SP_1 \cap T$ can be chosen, where $y_1 \not\in SZ(G)$. The element $y_1$ must be some $y_1 = sx_1$ for some non-identity elements $s \in S$ and $x_1 \in P_1$, where $x_1 \not\in Z(G)$. The commutator identity \ref{EqCommId} yields $\{1\} = [y_1, T] = [sx_1, T] = [s, T][x_1, T]$, which implies that $[s, T] = \{1\}$, where $S$ is abelian and $[s, S] = \{1\}$. Then $s$ centralises $S$ and $T$, which generate $G$, and so $s \in Z(G)$, a contradiction. Then $|SP_1 \cap T| = 1$, which is also a contradiction.
\end{proof}

\section{Groups of Nilpotency Class $2$}\label{SecGroupsOfTypeI}

The subgroup TPP ratio data for groups of nilpotency class $2$ \citep[subset of Tables 1-4]{HM} suggests an inexact (or strict) upper bound of $\sqrt{|G:Z(G)|}$. To minimise duplication Table 1 below summarises the relevant data for all groups of nilpotency class $2$ where $\rho_0$ is known that are \emph{not} $p$-groups with cyclic commutator subgroups of order $p$ (Tables \ref{tbl:TblRho0GroupsOfTypeIIa} and \ref{tbl:TblRho0GroupsOfTypeIIb}), with the group-theoretic data obtained from the GAP computer algebra system \citep{GAP}.

\begin{table}[H]
  \begin{center}
    \caption{Known subgroup TPP ratio data ($\rho_0$) for groups of nilpotency class $2$ which do not have a cyclic commutator subgroups of order $p$}
    \label{tbl:TblRho0GroupsOfTypeI}
    \begin{tabular}{lllllll}
      $\bm{|G|}$ & \multicolumn{1}{l}{\textbf{GAP ID}} & \multicolumn{1}{l}{\textbf{Structure Description}} & \multicolumn{1}{l}{$\bm{|Z(G)|}$} & \multicolumn{1}{l}{$\bm{\sqrt{|G:Z(G)|}}$} & \multicolumn{1}{l}{$\bm{\textbf{cd}(G)}$} & \multicolumn{1}{l}{$\bm{\rho_0}$} \\
      \toprule
        \multirow{2}{*}{$24$} &
            [24, 10] & $C_3 \times D_8$ & $6$ & $2$ & $\{1, 2\}$ & $1$ \\
          & [24, 11] & $C_3 \times Q_8$ & $6$ & $2$ & $\{1, 2\}$ & $1$ \\
        \hline
        \multirow{9}{*}{$32$} & 
            [32, 28] & $(C_4 \times C_2 \times C_2) \rtimes C_2$ & $4$ & $2\sqrt{2}$ & $\{1, 2\}$ & $1$ \\
        & [32, 29] & $(C_2 \times Q_8) \rtimes C_2$ & $4$ & $2\sqrt{2}$ & $\{1, 2\}$ & $1$ \\
        & [32, 30] & $(C_4 \times C_2 \times C_2) \rtimes C_2$ & $4$ & $2\sqrt{2}$ & $\{1, 2\}$  & $1$ \\
        & [32, 31] & $(C_4 \times C_4) \rtimes C_2$ & $4$ & $2\sqrt{2}$ & $\{1, 2\}$ & $1$ \\
        & [32, 32] & $(C_2 \times C_2) . (C_2 \times C_2 \times C_2)$ & $4$ & $2\sqrt{2}$ & $\{1, 2\}$ & $1$ \\
        & [32, 33] & $(C_4 \times C_4) \rtimes C_2$ & $4$ & $2\sqrt{2}$ & $\{1, 2\}$ & $1$ \\
        & [32, 34] & $(C_4 \times C_4) \rtimes C_2$ & $4$ & $2\sqrt{2}$ & $\{1, 2\}$ & $1$ \\
        & [32, 35] & $C_4 \rtimes Q_8$ & $4$ & $2\sqrt{2}$ & $\{1, 2\}$ & $1$ \\
        \hline
        \multirow{1}{*}{$64$} & 
            [64, 226] & $D_8 \rtimes D_8$ & $4$ & $4$ & $\{1, 2, 4\}$ & $2$ \\
        \hline
        \multirow{5}{*}{$128$} & 
            [128, 1135] & $(C_2 \times C_2 \times C_2 \times D_8) \rtimes C_2$ & $8$ & $4$ & $\{1, 2, 4\}$ & $2$ \\
         & [128, 1142] & $(C_2 \times C_2 \times ((C_4 \times C_2) \rtimes C_2)$ & $8$ & $4$ & $\{1, 2, 4\}$ & $2$ \\
         & [128, 1165] & $(C_2 \times C_2 \times C_2 \times D_8) \rtimes C_2$ & $8$ & $4$ & $\{1, 2, 4\}$ & $2$ \\
         & [128, 2194] & $C_2 \times D_8 \times D_8$ & $8$ & $4$ & $\{1, 2, 4\}$ & $2$ \\
         & [128, 2213] & $((C_4 \times D_8) \rtimes C_2) \rtimes C_2$ & $8$ & $4$ & $\{1, 2, 4\}$ & $2$ \\
    \end{tabular}
  \end{center}
\end{table}

The proof is given below.

\begin{theorem}\label{ThmNilpClassTwo} If $G$ is a group of nilpotency class $2$ then
\begin{equation}\label{SubTPPNilpClassTwo}
\rho_0(G) < \sqrt{|G: Z(G)|}.
\end{equation}
\end{theorem}
\begin{proof}\textup{(1)} The proof will be by induction on $|G|$ where $G$ is a group of nilpotency class $2$.

\vspace{1mm}

(Base case) The smallest groups of nilpotency class $2$ are the extraspecial groups of order $8 = 2^{1 + 2}$, namely, the dihedral group $D_8$ (GAP ID [8, 3]) and the quaternion group $Q_8$ (GAP ID [8, 4]): these contain cyclic (and normal) subgroups of index $2$ and, by \citep[Corollary 4.3]{Mur}, satisfy $\rho_0(G) = 1 < \sqrt{|G:Z(G)|} = \sqrt{2^2} = 2$.  Note that by Proposition \ref{PropRho0p4} $\rho_0 = 1$ is also true for all extraspecial groups of order $p^3$ for all primes $p$.

\vspace{1mm}

(Inductive hypothesis) Suppose (\ref{SubTPPNilpClassTwo}) is true for all groups of nilpotency class $2$ of order $ < |G|$, where $G$ is of nilpotency class $2$ and $|G| > 8$.

\vspace{1mm}

Let $(S, T, U)$ be a maximal non-trivial subgroup TPP triple of $G$, that is, $\frac{|S||T||U|}{|G|} = \rho_0(G) > 1$. Choosing any member, $S$, say, let $S_0 = S \cap Z(G)$, where $S_0$ is normal in $G$ (it is a subgroup of central elements of $G$). There are two cases: \textup{(i)} $S_0 = \{1\}$ and \textup{(ii)} $S_0 > \{1\}$, where in both cases $S_0 < S, Z(G)$.

\vspace{1mm}

Consider \textup{(i)}, $S_0 = \{1\}$. By Proposition \ref{PropTPPPermSubsets} and Observation \ref{ObsNilpClassTwo} $S$ is an abelian non-normal subgroup of $G$, and the product $SZ(G)$ is an abelian normal subgroup of $G$. Let $H = SZ(G)T$, where $H$ is normal in $G$. There are two subcases: \textup{(i.a)} $H < G$ and \textup{(i.b)} $H = G$.

\vspace{1mm}

Consider \textup{(i.a)}, $H < G$. Letting $U_0 = U \cap H$ and $U_1 = U \cap G\backslash H$ and applying Proposition \ref{PropTPPSplit} $\rho_0(G) = \frac{|S||T||U|}{|G|} \leq \rho_0(H)$. By Proposition \ref{PropTPPPermSubsets} $S$ and $T$ don't centralise each other, and there is a normal series for $H$ given by $\{1\} < H' = [SZ(G)T, SZ(G)T] = [S, T] \leq G' \leq Z(G) \leq Z(H) < H$, which yields a central series $\{1\} < H' \leq Z(H) < H$ of length $2$ and shows that $H$ is of nilpotency class $2$. The inductive hypothesis yields $\rho_0(H) < \sqrt{|H:Z(H)|}$, and the conditions $|H| < |G|$ and $|Z(H)| \geq |Z(G)|$ imply $\rho_0(G) \leq \rho_0(H) < \sqrt{|H:Z(H)|} \leq \sqrt{|G:Z(G)|}$.

\vspace{1mm}

Now consider \textup{(i.b)} $H = SZ(G)T = STZ(G) = \langle S, T, Z(G) \rangle = G$. Let $N = SZ(G) \cap T$. Then $N$ is normal in $G$ as it is normal in $SZ(G)$ and $T$, which generate $G$, and also $N \subseteq T$. Supposing $N > \{1\}$, by Lemma \ref{LemNeuQuoSubTPP}, the quotient $G/N$, which is a group of order $< |G|$, realises the subgroup TPP triple $(T/N,SN/N, UN/N)$ of type $(|T|/|N|,|S|,|U|)$ and $\rho_0(G/N) \geq \rho_0(G)$, and by the inductive hypothesis and Lemma \ref{LemQuoNilpClassTwo} $\rho_0(G) \leq \rho_0(G/N) < \sqrt{|G/N|:|Z(G/N)|} \leq \sqrt{|G:Z(G)|}$. Suppose $N = \{1\}$. Then $SZ(G)T = G$ is an (internal) semidirect product $SZ(G) \rtimes T$, and $G/SZ(G) \cong T$ and $|S||T||Z(G)| = |G|$. Permuting $T$ and $U$, which corresponds to the permuted TPP triple $(S, U, T)$, it can be shown that $G/SZ(G) \cong U$, which implies that $T \cong U$ and $|T| = |U|$. By also considering the permuted triples $(T, S, U)$ and $(T, U, S)$ it can be shown that $|S| = |U|$, and by transitivity $|S| = |T| = |U|$. Finally, from $|G| = |S||T||Z(G)| = |S|^2|Z(G)|$ it follows that $|S| = |T| = |U| = \sqrt{|G:Z(G)|}$ and $\rho_0(G) = \frac{|S||T||U|}{|G|} = \frac{|S|^3}{|G|} = \frac{\sqrt{|G:Z(G)|}}{|Z(G)|} < \sqrt{|G:Z(G)|}$.

\vspace{1mm}

Consider \textup{(ii)} $S_0 = S \cap Z(G) > \{1\}$, where $S_0 < S, Z(G)$. Then $G/S_0$ is a group of order $< |G|$ and nilpotency class $2$ as there is a central series of length $2$ given by $\{1\} < (G/S_0)' = [G/S_0, G/S_0] = G'S_0/S_0 \leq Z(G)S_0/S_0 < G/S_0$. The inductive hypothesis and Lemma \ref{LemQuoNilpClassTwo} yield $\rho_0(G/S_0) < \sqrt{|G/S_0|:|Z(G/S_0)|} \leq \sqrt{|G:Z(G)|}$. Lemma \ref{LemNeuQuoSubTPP} implies that $G/S_0$ realises the subgroup TPP triple $(S/S_0, TS_0/S_0, US_0/S_0)$ of type $(|S|/|S_0|,|T|,|U|)$ and $\rho_0(G) = \frac{|S||T||U|/|S_0|}{|G|/|S_0|} \leq \rho_0(G/S_0) < \sqrt{|G:Z(G)|}$. By permutation invariance the same conclusion follows from supposing $T_0 = T \cap Z(G) > \{1\}$ or $U_0 = U \cap Z(G) > \{1\}$.

\end{proof}

The result may be compared with a fact from the character theory of finite groups that if $d$ is the degree of an (irreducible) character of a group $G$ then $d \leq \sqrt{|G:Z(G)|}$.

\section{$p$-Groups with a Cyclic Commutator Subgroup of Order $p$}\label{SecGroupsOfTypeII}

The $p$-groups with a cyclic commutator subgroup of order $p$ are, as proved in Lemma \ref{LemCycNormSubOrderp}, of nilpotency class $2$. Any member $G$ of this class of order, $p^n$ for $n \geq 3$, say, contains a centre of order at least $p$, and by Theorem \ref{ThmNilpClassTwo} there is a variable bound in $p$ and $n$ for $\rho_0(G)$ given by

\begin{align}\rho_0(G) < p^{\frac{(n - 1)}{2}}.
\end{align}

A better and sharp bound of $p$ is suggested by the very limited subgroup TPP data \citep[Tables 1 \& 2]{HM} for this class of groups for small $p$, namely, $p=2, 3$, as reproduced below in Table 2, showing the (qualifying) $2$-groups of order $8$, $16$, and $32$, and Table 3, showing the $3$-groups of order $27$.

\begin{table}[H]
  \begin{center}
    \caption{Known subgroup TPP ratio data ($\rho_0$) for small $2$-groups with a cyclic commutator subgroup of order $2$}
    \label{tbl:TblRho0GroupsOfTypeIIa}
    \begin{tabular}{llllllll}
      $\bm{|G|}$ & \multicolumn{1}{l}{\textbf{GAP ID}} & \multicolumn{1}{l}{\textbf{Structure Description}} & \multicolumn{1}{l}{$\bm{|Z(G)|}$} & \multicolumn{1}{l}{$\bm{\sqrt{|G:Z(G)|}}$} & \multicolumn{1}{l}{$\bm{\textbf{cd}(G)}$} & \multicolumn{1}{l}{$\bm{\rho_0}$} \\
      \toprule
        \multirow{2}{*}{$8$} &
            [8, 3] & $D_8$ & $2$ & $2$ & $\{1, 2\}$ & $1$ \\
          & [8, 4] & $Q_8$ & $2$ & $2$ & $\{1, 2\}$ & $1$ \\
        \hline
        \multirow{6}{*}{$16$} & 
            [16, 3] & $(C_4 \times C_2) \rtimes C_2$ & $4$ & $2$ & $\{1, 2\}$ & $1$ \\
          & [16, 4] & $C_4 \times C_4$ & $4$ & $2$ & $\{1, 2\}$ & $1$ \\
          & [16, 6] & $C_8 \rtimes C_2$ & $4$ & $2$ & $\{1, 2\}$ & $1$ \\
          & [16, 11] & $C_2 \times D_8$ & $4$ & $2$ & $\{1, 2\}$ & $1$ \\
          & [16, 12] & $C_2 \times Q_8$ & $4$ & $2$ & $\{1, 2\}$ & $1$ \\
          & [16, 13] & $(C_4 \times C_2) \rtimes C_2$ & $4$ & $2$ & $\{1, 2\}$ & $1$ \\
        \hline
        \multirow{5}{*}{$32$} &
            [32, 2] & $(C_4 \times C_2) \rtimes C_4$ & $8$ & $2$ & $\{1, 2\}$ & $1$ \\
          & [32, 4] & $C_8 \rtimes C_4$ & $8$ & $2$ & $\{1, 2\}$ & $1$ \\
          & [32, 5] & $(C_8 \times C_2) \rtimes C_2$ & $8$ & $2$ & $\{1, 2\}$ & $1$ \\
          & [32, 12] & $C_4 \times C_8$ & $8$ & $2$ & $\{1, 2\}$ & $1$ \\
          & [32, 17] & $C_{16} \rtimes C_2$ & $8$ & $2$ & $\{1, 2\}$ & $1$ \\
          & [32, 22] & $C_2 \times ((C_4 \times C_2) \rtimes C_2)$ & $8$ & $2$ & $\{1, 2\}$ & $1$ \\
          & [32, 23] & $C_2 \times (C_4 \rtimes C_4)$ & $8$ & $2$ & $\{1, 2\}$ & $1$ \\
          & [32, 24] & $(C_4 \times C_4) \rtimes C_2$ & $8$ & $2$ & $\{1, 2\}$ & $1$ \\
          & [32, 25] & $C_4 \times D_8$ & $8$ & $2$ & $\{1, 2\}$ & $1$ \\
          & [32, 26] & $C_4 \times Q_8$ & $8$ & $2$ & $\{1, 2\}$ & $1$ \\
          & [32, 37] & $C_2 \times (C_8 \rtimes C_2)$ & $8$ & $2$ & $\{1, 2\}$ & $1$ \\
          & [32, 38] & $(C_8 \times C_2) \rtimes C_2$ & $8$ & $2$ & $\{1, 2\}$ & $1$ \\
          & [32, 46] & $C_2 \times C_2 \times D_8$ & $8$ & $2$ & $\{1, 2\}$ & $1$ \\
          & [32, 47] & $C_2 \times C_2 \times Q_8$ & $8$ & $2$ & $\{1, 2\}$ & $1$ \\
          & [32, 48] & $C_2 \times ((C_4 \times C_2) \times C_2)$ & $8$ & $2$ & $\{1, 2\}$ & $1$ \\
          & [32, 49] & $(C_2 \times C_2 \times C_2) \rtimes (C_2 \times C_2)$ & $2$ & $4$ & $\{1, 4\}$ & $2$ \\
          & [32, 50] & $(C_2 \times Q_8) \rtimes C_2$ & $2$ & $4$ & $\{1, 2\}$ & $1$ \\
    \end{tabular}
  \end{center}
\end{table}

\begin{table}[H]
  \begin{center}
    \caption{Known subgroup TPP ratio data ($\rho_0$) for small $3$-groups with a cyclic commutator subgroup of order $3$}
    \label{tbl:TblRho0GroupsOfTypeIIb}
    \begin{tabular}{llllllll}
      $\bm{|G|}$ & \multicolumn{1}{l}{\textbf{GAP ID}} & \multicolumn{1}{l}{\textbf{Structure Description}} & \multicolumn{1}{l}{$\bm{|Z(G)|}$} & \multicolumn{1}{l}{$\bm{\sqrt{|G:Z(G)|}}$} & \multicolumn{1}{l}{$\bm{\textbf{cd}(G)}$} & \multicolumn{1}{l}{$\bm{\rho_0}$} \\
      \toprule
        \multirow{2}{*}{$27$} &
            [27, 3] & $(C_3 \times C_3) \rtimes C_3$ & $3$ & $3$ & $\{1, 3\}$ & $1$ \\
          & [27, 4] & $C_9 \rtimes C_3$ & $3$ & $3$ &  $\{1, 3\}$ &$1$ \\
    \end{tabular}
  \end{center}
\end{table}

(There is some overlap in these tables with the $p$-groups of nilpotency class $2$ satisfying the large centre assumption $p^2 \leq |G:Z(G)| \leq p^3$ or satisfying the small character degrees assumption $\text{cd}(G) = \{1, p\}$.) The data show that $\rho_0 = 1$, except for a single $2$-group of order $32$, which is an extraspecial group (GAP ID [32, 49]) that does reach the upper bound of $2$. This suggests a bound of $\rho_0 \leq p$ for this class of groups, which is proved below.

\begin{theorem}\label{ThmRho0pGroupCycComm} If $G$ is a nonabelian $p$-group with a cyclic commutator subgroup $G'$ of order $p$ then
\begin{equation}\label{EqnRho0pGroupCycComm} \rho_0(G) \leq p.
\end{equation}
\end{theorem}
\begin{proof} The proof will be by induction on $|G| = p^n$, $n \geq 3$, where $G$ is a nonabelian $p$-group with a cyclic commutator subgroup $G'$ of order $p$, and $p$ is an arbitrary fixed prime. Recall that, as shown in Lemma \ref{LemCycNormSubOrderp}, they satisfy \ref{NilpClassTwo}. The base case is extraspecial groups of order $8$, which have already been treated in Theorem \ref{ThmNilpClassTwo} and shown to have $\rho_0 = 1$. Note that by Proposition \ref{PropRho0p4} $\rho_0 = 1$ is also true for all extraspecial groups of order $p^3$ for all primes $p$.

(Inductive hypothesis) Suppose \ref{EqnRho0pGroupCycComm} is true for all $p$-groups of order $< p^n$, for $n \geq 5$, with cyclic commutator subgroups of order $p$.

\vspace{1mm}

The approach taken here will be almost identical to that used in Theorem \ref{ThmNilpClassTwo}. Let $(S, T, U)$ be a maximal non-trivial subgroup TPP triple of $G$ such that $\frac{|S||T||U|}{|G|} = \rho_0(G) > 1$. Let $S_0 = S \cap Z(G)$. Then $S_0$ is normal in $G$ as it is central in $G$. There are two cases to consider: \textup{(i)} $S_0 = \{1\}$ and \textup{(ii)} $S_0 > \{1\}$, where in both cases $S_0 < S, Z(G)$. The easiest case is \textup{(ii)} $S_0 > \{1\}$. Then $|S_0| \geq p$ and there is a quotient $G/S_0$ of order $< p^n$ which, by Lemma \ref{LemNeuQuoSubTPP} and the inductive hypothesis, s{}atisfies $\rho_0(G) \leq \rho_0(G/S_0) \leq p$.

\vspace{1mm}

Consider case \textup{(i)}, $S_0 = \{1\}$. As in Theorem \ref{ThmNilpClassTwo} $T \cap Z(G) = \{1\}$ and $U \cap Z(G) = \{1\}$ may also be supposed. Let $H = SZ(G)T = STZ(G)$. Then, by Observation \ref{ObsNilpClassTwo}, $S$ (and also $T$ and $U$) is abelian and $H$ is normal in $G$ . There are two subcases \textup{(i.a)} $H < G$ and \textup{(i.b)} $H = G$. Consider \textup{(i.a)}, $H < G$. Then $H$ is of order $< p^n$, and of nilpotency class $2$ as there is a central series of length $2$ given by $\{1\} < H' = [H, H] \leq Z(H) < H$, where $H' = [SZ(G)T, SZ(G)T] = [S, T]$, and by the inductive hypothesis and Proposition \ref{PropTPPSplit} $\rho_0(G) \leq \rho_0(H) \leq p$.

\vspace{1mm}

Now consider \textup{(i.b)} $H = SZ(G)T = \langle S, T, Z(G) \rangle = G$. It will be shown, by constructing in parallel another normal subgroup $K = SG'T = STG'$ of $G$, and comparing it with $H$ while using the assumption that $G'$ is cyclic of order $p$, that $G$ must be extraspecial. First, the possibility that $K < G$ can be dealt with using the same argument above for the subcase \textup{(i.a)}, with $K$ replacing $H$ and $G'$ replacing $Z(G)$, which implies that $\rho_0(G) \leq \rho_0(K) \leq p$. So consider $K = SG'T = G = SZ(G)T$, where $SG'$ is normal in $G$. As shown in Theorem \ref{ThmNilpClassTwo} any non-trivial intersection $N = SZ(G) \cap T > \{1\}$ can be used, together with the inductive hypothesis, to construct a quotient $G/N$ such that $\rho_0(G) \leq \rho_0(G/N) \leq p$. Let $SZ(G) \cap T = \{1\}$ Then $G$ is a semidirect product $G = SZ(G) \rtimes T$. Now, as $S \cap G' = \{1\}$ and $SG' \cap T = \{1\}$, the condition $K = SG'T = SZ(G)T$ implies that $G$ is also the semidirect product $G = SG' \rtimes T$, which implies that $TG' \cong TZ(G)$ and $G' \cong Z(G) \cong \mathbb{Z}_p$. The Frattini subgroup, which is the intersection of all the maximal subgroups of $G$ and consists of all the non-generating elements of $G$, is the product $\Phi(G) = G^pG'$, where $G^p$ is the subgroup generated by all the $p$-th powers of elements of $G$, and $G/\Phi(G)$ is elementary abelian. In the case being considered $G' \cong \mathbb{Z}_p$, so all commutators $[x, y] \in G'$ are of order $p$, and \ref{EqCommId2} yields $[g^p, x] = [g, p]^p = 1$ for any $g^p \in G^p$ and any $x \in G$. That is, $G^p$ centralises $G$ and so $G^p \leq Z(G) = G'$, implying that $G' = Z(G) = \Phi(G)$ and that $G$ is extraspecial of order $p^{1 + 2m}$ for some $m \geq 2$ such that $n = 2m + 1$.

\vspace{1mm}

As shown in Theorem \ref{ThmNilpClassTwo}, by permuting $S, T, U$, there are also semidirect products $G = SZ(G) \rtimes U = TZ(G) \rtimes S =  TZ(G) \rtimes U$, which imply that $S \cong T \cong U$ and $|S| = |T| = |U| = \sqrt{|G:Z(G)|} = p^m$. Also $S, T, U$ are all abelian (from $S \cap G' = T \cap G' = U \cap G' = \{1\}$). Writing $G$ as $G = S\Phi(G)T = \langle S, T, \Phi(G) \rangle$ the Frattini subgroup $\Phi(G)$ can be removed, which implies that $\langle S, T \rangle = G$. But this is directly contradicted by Lemma \ref{LemExspSub}.
\end{proof}

This leads to an immediate corollary for extraspecial groups, which have the property that $G'$ is cyclic of order $p$ such that $G' = Z(G) = \Phi(G)$.

\begin{corollary}\label{CorrRho0pGroupCycComm} If $G$ is an extraspecial group of order $p^{1 + 2n}$, for $n \geq 1$, then $\rho_0(G) \leq p$.
\end{corollary}

\section{$p$-Groups Of Nilpotency Class $2$ with a Large Centre}\label{SecGroupsOfTypeIII}

The groups considered here are $p$-group of nilpotency class $2$ that satisfy $p^2 \leq |G:Z(G)| \leq p^3$. The available subgroup TPP ratio data for these groups can be found as a subset of Table \ref{tbl:TblRho0GroupsOfTypeI}, and Tables \ref{tbl:TblRho0GroupsOfTypeIIa} and \ref{tbl:TblRho0GroupsOfTypeIIb}, where $|G:Z(G)| = p^3$. Although Theorem \ref{ThmNilpClassTwo} implies that $\rho_0 \leq p$ here, a stronger result of $\rho_0 = 1$ follows from Proposition \ref{PropAbCentQuop2p3}.

\begin{theorem}\label{ThmRho0pGroupLargeCentre} If $G$ is a $p$-group of nilpotency class $2$ such that $p^2 \leq |G:Z(G)| \leq p^3$ then $\rho_0(G) = 1$.
\end{theorem}
\begin{proof} Let $G$ be as given. By Theorem \ref{NilpClassTwo} $\rho_0(G) \leq p$, which means that the size $|S||T||U|$ of any subgroup TPP triple of $G$ satisfies $|S||T||U| \leq p|G| = p^{n + 1}$. However, by Proposition \ref{PropAbCentQuop2p3} $G$ contains abelian (maximal, normal) subgroups of index $p$, and by Corollary \citep[Corollary 4.3]{Mur} $\rho_0(G) = 1$.
\end{proof}

\section{$p$-Groups of Nilpotency Class $2$ with Small Character Degrees}\label{SecGroupsOfTypeIV}

As proved in \citep[Theorem 4.1, Corollary 4.3]{Mur} and used in Theorem \ref{ThmRho0pGroupLargeCentre} groups with large abelian normal subgroups of prime index $p$ tend to have a correspondingly small subgroup TPP ratio $\rho_0$ not exceeding $p$. There is an interesting fact from the character theory of finite groups, that the degree of an irreducible character of a group with an abelian (though, not necessarily normal) subgroup of index $n$, where $n$ is not necessarily prime, cannot exceed $n$. The "smallness" of character degrees of a group also seems to play a role in bounding the value of $\rho_0$, at least indirectly by reflecting the existence of large abelian subgroups. The available subgroup TPP ratio data for $p$-groups of nilpotency class $2$ with $\text{cd}(G) = \{1, p\}$, which can be found as a subset of Table \ref{tbl:TblRho0GroupsOfTypeI}, and Tables \ref{tbl:TblRho0GroupsOfTypeIIa} and \ref{tbl:TblRho0GroupsOfTypeIIb}, certainly suggest that such groups satisfy $\rho_0 = 1$, as there is no counterexample, and this is proved in the following result.

\begin{theorem} If $G$ is a $p$-group of nilpotency class $2$ with character degree set $\text{cd}(G) = \{1, p\}$ then $\rho_0(G) = 1$.
\end{theorem}
\begin{proof}By \citep[Theorem C 4.8]{IP} any (nonabelian) $p$-group with $\text{cd}(G) = \{1, p\}$ either contains an abelian (maximal, normal) subgroup of index $p$, or has a centre of index $p^3$. In the first case $\rho_0 = 1$ by \citep[Corollary 4.3]{Mur}. In the second case $G/Z(G)$ is abelian of order $p^3$ and by Theorem \ref{ThmRho0pGroupLargeCentre} $\rho_0(G) = 1$.
\end{proof}

\hypertarget{bib}{}
\begin{thebibliography}{}
\bibitem{AGW}J. D. Adler, and M. Garlow and E. R. Wheland, 'Groups of order $p^4$ made less difficult', arXiv:1611.00461, 2016. \url{https://doi.org/10.48550/arXiv.1611.00461}
\bibitem{AB}J. L. Alperin and R. B. Bell, \emph{Groups and Representations}, Springer, 1991.
\bibitem{Ber}Y. Berkovich, \emph{Groups of Prime Power Order}, Volume 1, De Gruyter, 2008.
\bibitem{CU}H. Cohn and C. Umans, 'A group-theoretic approach to fast matrix multiplication', in: \emph{Proceedings of the 44th Annual Symposium on Foundations of Computer Science} (Cambridge, Massachusetts 2003), IEEE Computer Society (2003), 438-449. \url{https://dl.acm.org/doi/10.5555/946243.946301}
\bibitem{Cra}, D. A. Craven, The Theory of $p$-Groups, Lecture notes, 2008. \url{https://web.mat.bham.ac.uk/D.A.Craven/docs/lectures/pgroups.pdf}
\bibitem{IP}I. M. Isaacs and D. S. Passman, 'A characterization of groups in terms of the degrees of their characters', \emph{Pacific Journal of Mathematics}, 15(3): 877-903, 1965. \url{http://dx.doi.org/10.2140/pjm.1965.15.877}
\bibitem{GAP}The GAP Group, GAP - Groups, Algorithms, and Programming, Version 4.15.1, 2025. \url{https://www.gap-system.org}
\bibitem{HM}I. Hedtke and S. Murthy, 'Search and test algorithms for triple product property triples', \emph{Groups - Complexity - Cryptology}, 4(1): 111-133, 2012. \url{https://doi.org/10.1515/gcc-2012-0006}
\bibitem{Mur}S. R. Murthy, 'A note on the triple product property for finite groups with abelian normal subgroups of prime index', arXiv:2512.16730, 2025. \url{https://doi.org/10.48550/arXiv.2512.16730}
\bibitem{Neu2011}P. M. Neumann, 'A note on the triple product property for subsets of finite groups', \emph{LMS J. Computation and Mathematics}, 14:232-237, 2011. \url{https://doi.org/10.1112/S1461157010000288}
\bibitem{Neu2012}P. M. Neumann, Private communication, 2012.
\bibitem{Wil}M. W. Wild, 'The groups of order sixteen made easy', \emph{The American Mathematical Monthly}, 112(1): 20-31, 2005. \url{https://doi.org/10.2307/30037381}
\end{thebibliography}

\end{document}

%-----------------------------------------------------------------------
% End of amsart-template.tex
%-----------------------------------------------------------------------
